%----------------------------------------------------------------------------------------
% NOTATIE EN FIGUREN VAN DE CURSUS TELPROBLEMEN (2de graad)
% Staat in de preambule en niet in een hoofdstuk: ook de lijsten met oplossingen gebruiken dit.
% Vertrokken van notatie.tex van Integralen; wat hier niet nodig is, is geschrapt.
%----------------------------------------------------------------------------------------
% Bewaarde grafiek, geplaatst door maak-pdf.py: bestand en pagina in figuren-cache/
\newcommand{\grafiekbestand}[2]{\includegraphics[page=#2,trim=3bp 3bp 3bp 3bp]{figuren-cache/#1}}

% Zoals stappen, maar met een bredere kolom voor de berekening: voor lange formules
\newenvironment{stappenbreed}{%
  \par\renewcommand{\arraystretch}{1.25}%
  \noindent\begin{tabular}{@{}>{$\displaystyle}p{9.8cm}<{$}@{\hspace{4mm}}>{\RaggedRight\small\itshape\arraybackslash}p{4.9cm}@{}}%
}{\end{tabular}\par}

% Zoals stappen, met eigen kolombreedtes: \begin{stappenvrij}{6.2cm}{3.8cm}. Voor naast een figuur.
\newenvironment{stappenvrij}[2]{%
  \par\renewcommand{\arraystretch}{1.25}%
  \noindent\begin{tabular}{@{}>{$\displaystyle}p{#1}<{$}@{\hspace{3mm}}>{\RaggedRight\small\itshape\arraybackslash}p{#2}@{}}%
}{\end{tabular}\par}

% Tekst of berekening links, figuur rechts: \naast{links}{rechts}
\newsavebox{\naastbox}
\newcommand{\naast}[3][]{%
  \par\noindent\sbox{\naastbox}{#3}%
  \begin{minipage}[t]{\dimexpr\linewidth-\wd\naastbox-3mm\relax}\vspace{0pt}#2\end{minipage}\hfill
  \begin{minipage}[t]{\wd\naastbox}\vspace{0pt}\usebox{\naastbox}\end{minipage}\par\addvspace{1mm}}
% Een rij figuren, gecentreerd
\newcommand{\figuren}[1]{\par\addvspace{1mm}{\centering #1\par}\addvspace{1.5mm}}
% Opgave met een figuur rechts: \begin{metgrafiek}{tikz} opgave en deelvragen \end{metgrafiek}
\newsavebox{\grafiekbox}
\newenvironment{metgrafiek}[2][]{%
  \sbox{\grafiekbox}{#2}%
  \begin{minipage}[t]{\dimexpr\linewidth-\wd\grafiekbox-3mm\relax}%
}{%
  \end{minipage}\hfill
  \begin{minipage}[t]{\wd\grafiekbox}\vspace{-2mm}\usebox{\grafiekbox}\end{minipage}\par}

%----------------------------------------------------------------------------------------
% BOOMDIAGRAM, binnen een tikzpicture, met de syntaxis child{node{..}} van TikZ:
%   \begin{tikzpicture}[boom=18mm, level 1/.style={sibling distance=16mm}, level 2/.style={sibling distance=8mm}]
%   \node[wortel] {} child{node{K} child{node{K}} child{node{M}}} child{node{M} ...};
% boom=<afstand tussen de niveaus>; de takken groeien naar rechts, in de kleur van het hoofdstuk.
% uitkomst: een grijs label rechts van een blad, met de uitkomst, bv. node[uitkomst=KKM]{M}.
%----------------------------------------------------------------------------------------
\tikzset{
  boom/.style={grow'=right,level distance=#1,
    every node/.style={font=\small,inner sep=1.5pt},
    edge from parent/.style={draw=accent,thick},
    edge from parent path={(\tikzparentnode.east) -- (\tikzchildnode.west)}},
  boom/.default=16mm,
  wortel/.style={circle,fill=accent,inner sep=1.6pt},
  knoop/.style={circle,fill=accent,inner sep=1.2pt},
  uitkomst/.style={label={[font=\scriptsize,text=grijs,inner sep=1pt]right:#1}},
  tak/.style={accent,thick},
  blad/.style={font=\small,inner sep=1.5pt},
  plaats/.style={font=\small,inner sep=1.5pt,fill=white},
  verz/.style={draw=accent,thick,fill=accent,fill opacity=.10,text opacity=1},
  stad/.style={circle,fill=accent,inner sep=1.8pt},
}

%----------------------------------------------------------------------------------------
% VENNDIAGRAM MET TWEE VERZAMELINGEN, binnen een tikzpicture.
% \venntwee{naam A}{naam B}: de omgevingsverzameling U en twee cirkels. Daarna zijn er coördinaten
% (vA) alleen in A, (vAB) in beide, (vB) alleen in B, (vU) buiten de cirkels.
% \vennaantal{alleen A}{beide}{alleen B}{buiten}: de aantallen in de vier stukken (leeg laten mag).
% Een stuk kleuren vóór \venntwee: \fill[accent!35] \cirkelA; of met \clip \cirkelA; \fill \cirkelB;
% \vennU geeft de rechthoek, voor het complement: \fill[accent!35,even odd rule] \vennU \cirkelA;
%----------------------------------------------------------------------------------------
\newcommand{\cirkelA}{(0,0) circle (.95)}
\newcommand{\cirkelB}{(1.1,0) circle (.95)}
\newcommand{\vennU}{(-1.45,-1.25) rectangle (3.05,1.3)}
\newcommand{\venntwee}[2]{%
  \draw[grijs] \vennU;
  \node[font=\scriptsize,text=grijs,anchor=south west,inner sep=2pt] at (-1.45,-1.25) {$U$};
  \draw[verz] \cirkelA; \draw[verz] \cirkelB;
  \node[font=\footnotesize,text=accent,anchor=south east,inner sep=1pt] at (-.5,.8) {#1};
  \node[font=\footnotesize,text=accent,anchor=south west,inner sep=1pt] at (1.6,.8) {#2};
  \coordinate (vA) at (-.4,0); \coordinate (vAB) at (.55,0); \coordinate (vB) at (1.5,0);
  \coordinate (vU) at (2.65,-.95);}
\newcommand{\vennaantal}[4]{%
  \node at (vA) {#1}; \node at (vAB) {#2}; \node at (vB) {#3}; \node at (vU) {#4};}

% VENNDIAGRAM MET DRIE VERZAMELINGEN: \venndrie{A}{B}{C}, met coördinaten (vA), (vB), (vC) (alleen in die
% verzameling), (vAB), (vAC), (vBC) (in precies die twee), (vABC) (in alle drie) en (vU) (buiten).
\newcommand{\cirkelP}{(0,0) circle (1)}
\newcommand{\cirkelQ}{(1.2,0) circle (1)}
\newcommand{\cirkelR}{(.6,-1.03) circle (1)}
\newcommand{\vennUdrie}{(-1.5,-2.35) rectangle (2.7,1.4)}
\newcommand{\venndrie}[3]{%
  \draw[grijs] \vennUdrie;
  \node[font=\scriptsize,text=grijs,anchor=south west,inner sep=2pt] at (-1.5,-2.35) {$U$};
  \draw[verz] \cirkelP; \draw[verz] \cirkelQ; \draw[verz] \cirkelR;
  \node[font=\footnotesize,text=accent,anchor=south east,inner sep=1pt] at (-.55,.82) {#1};
  \node[font=\footnotesize,text=accent,anchor=south west,inner sep=1pt] at (1.75,.82) {#2};
  \node[font=\footnotesize,text=accent,anchor=north,inner sep=1pt] at (.6,-2.06) {#3};
  \coordinate (vA) at (-.4,.2); \coordinate (vB) at (1.6,.2); \coordinate (vC) at (.6,-1.5);
  \coordinate (vAB) at (.6,.42); \coordinate (vAC) at (-.05,-.68); \coordinate (vBC) at (1.25,-.68);
  \coordinate (vABC) at (.6,-.35); \coordinate (vU) at (2.3,-2.05);}
\newcommand{\vennaantaldrie}[8]{%
  \node at (vA) {#1}; \node at (vB) {#2}; \node at (vC) {#3}; \node at (vAB) {#4};
  \node at (vAC) {#5}; \node at (vBC) {#6}; \node at (vABC) {#7}; \node at (vU) {#8};}
